\documentclass[11pt]{article}
\usepackage[margin=25mm]{geometry}
\usepackage{url}
\usepackage{enumitem}
\usepackage{xcolor}
\usepackage{hyperref}

\newcommand{\condition}[2]{\section*{Condition #1: #2}}

\begin{document}

\begin{center}
{\Large Response to Reviewers}\\[4pt]
Paper No. S-28-4-002\\
Original title: Camerala: A Privacy-Preserving Browser-Native Framework for In-the-Wild Physiological Sensing and Task Measurement
\end{center}

We thank the handling editor and the reviewers for their careful reading and constructive comments.
The manuscript was revised according to Reviewer 1's required Conditions 1--13.
The revision narrows unsupported claims while preserving the intended contribution: a local-first browser-native framework for remote psychological/psychophysiological task sensing.
We also thank Reviewer 2 for the positive evaluation.

\condition{1}{The title of the manuscript should be reconsidered}
\textbf{Response.}
Thank you for this comment.
The manuscript was revised to avoid implying that Camerala proposes a dedicated privacy-preserving algorithm.
Following the overall revision, the title was changed to a more conservative local-first system title for remote psychophysiological task sensing.

\textbf{Revision.}
The title no longer uses ``Privacy-Preserving'' or an unqualified ``In-the-Wild'' expression.
Privacy is now discussed as an architectural privacy advantage of avoiding raw-video transmission, not as a dedicated privacy-preserving method.
The new title also preserves the intended remote psychophysiological task-sensing use case.

\textbf{Key revised text.}
``Camerala: A Local-First Browser-Native Framework for Remote Psychophysiological Task Sensing''

\textbf{Location.}
Title page; Introduction; Contributions.

\condition{2}{The necessity of real-time processing is not sufficiently explained}
\textbf{Response.}
Thank you for this important clarification.
We revised the manuscript to explain that real-time processing is not claimed as the only technically possible architecture.
We acknowledge that local video recording followed by offline processing could preserve timestamps if implemented carefully.

\textbf{Revision.}
The revised manuscript clarifies that Camerala targets remote psychological and psychophysiological task experiments conducted in private spaces.
In such settings, retaining identifiable facial video for later analysis can itself be a privacy and acceptability burden.
Camerala therefore extracts face-image-derived records during task execution and stores only task events, exported feature fields, and auxiliary context fields, without raw facial video persistence.
The role of real-time browser-side processing is thus feature-only local logging and timestamp alignment within the same client-side runtime, not online task adaptation or a claim that physiological estimates are more accurate.

\textbf{Key revised text.}
``Local video recording followed by offline processing could preserve timestamps if implemented carefully. However, the intended use case of Camerala is a remote psychological or psychophysiological task experiment conducted in a private space, where retaining identifiable facial video for later analysis can itself be a privacy and acceptability burden.''

\textbf{Location.}
Introduction; System Architecture; Data Management and Local-First Persistence.

\condition{3}{The Research Questions are not defined clearly enough}
\textbf{Response.}
Thank you for this comment.
The Research Questions were rewritten to define the tested configuration and the observation criteria more concretely.

\textbf{Revision.}
We revised RQ1 to ask whether Camerala's task-synchronized face-derived records, including rPPG-oriented ROI values, satisfy the operational criteria used for descriptive analysis in the intended remote psychological/psychophysiological task use case.
The revised manuscript defines concrete operational criteria at the trial and exported-window levels: reaction times within 100--1500 ms for trial-level retention and \texttt{mean\_quality} $\geq$ 0.8 for exported-window-level landmark availability.
The completed deployment consisted of 600 trials in total.
After applying the RT criterion, 541/600 trials were retained for descriptive analysis.
All 299 exported windows had \texttt{mean\_quality} = 1.0 and therefore satisfied the operational landmark-availability criterion.

\textbf{Key revised text.}
``For RQ1, we evaluated the exported records using operational criteria at the trial and exported-window levels. At the trial level, trials were retained for descriptive analysis when reaction times fell within 100--1500 ms. At the exported-window level, landmark availability was evaluated using \texttt{mean\_quality}, the mean of the binary landmark-availability flag within each 10-s window; exported windows with \texttt{mean\_quality} $\geq$ 0.8 were treated as satisfying the operational landmark-availability criterion.''

\textbf{Location.}
Introduction, Research Questions.

\condition{4}{It is unclear what the ``between-environment differences'' in RQ2 refer to}
\textbf{Response.}
Thank you for this important clarification.
We retained the RQ2 wording and clarified its interpretation as a descriptive comparison of task-synchronized face-derived and sensor-derived exported records between laboratory and home deployment contexts.
The revised manuscript explicitly states that these differences may reflect environmental conditions, screen light, posture, camera geometry, task timing, the measurement pipeline, participant state, or their mixture, and are not interpreted as psychological or physiological effects.

\textbf{Revision.}
The Abstract, the explanatory text immediately following RQ2, Results, Discussion, and Conclusion were revised to state that the observed laboratory/home differences are deployment-context differences in exported records.
Because the present design did not include reference physiological sensors, multiple participants, or controlled separation of lighting, screen light, posture, camera geometry, task timing, measurement-pipeline factors, and participant state, the observed differences are not interpreted as psychological, cognitive, motivational, or physiological effects.
Future work requirements were added, including reference ECG or contact PPG, multiple participants, multiple devices, controlled lighting/camera/posture conditions, and manipulation checks.

\textbf{Key revised text.}
``RQ2 was retained as a question about between-environment differences, but it was treated as a descriptive question about deployment-context differences in exported records. In the revised interpretation, the phrase `physiological signals' refers conservatively to task-synchronized face-derived and sensor-derived exported records rather than validated physiological state measurements.''

``The present design was not able to determine whether laboratory/home differences reflected psychological state, physiological state, lighting, screen light, posture, camera geometry, the measurement pipeline, task timing, participant state, or their mixture. Therefore, these differences were not interpreted as psychological, cognitive, motivational, or physiological effects.''

\textbf{Location.}
Abstract; Introduction, Research Questions; Results; Discussion; Conclusion.

\condition{5}{The novelty and positioning of the study are insufficiently justified}
\textbf{Response.}
Thank you for identifying the relevant prior work.
We agree that browser-based, web-deployable, edge-based, and privacy-oriented rPPG systems have already been explored.
The Related Work section was expanded, and the contribution claim was repositioned accordingly.

\textbf{Revision.}
The Related Work section now includes all studies, systems, and public demonstrations explicitly identified by the reviewer: Gupta and Etemad, Ayesha et al., LabVanced rPPG, FacePhys-Demo, RhythmEdge, and Di Lernia et al.
LabVanced and FacePhys-Demo are described as public tools/demonstrations rather than peer-reviewed empirical validation studies.
We removed or substantially weakened contribution claims based on browser-based rPPG functionality, local rPPG processing, edge-oriented heart-rate estimation, or privacy preservation as an algorithm.
The revised manuscript instead positions Camerala as a system-integration contribution for remote psychological/psychophysiological task experiments.
Specifically, Camerala integrates browser-based task execution, face-derived feature extraction, timestamp alignment between task events and exported frame/window records, operational landmark-availability checking, and local CSV/ZIP persistence without raw facial video retention within one browser runtime.
This contribution is presented as a limited research prototype and implementation example, not as a new rPPG estimator or a validated physiological measurement system.

\textbf{Key revised text.}
``The contribution of Camerala is not browser-based rPPG functionality, local rPPG processing, edge-based heart-rate estimation, or privacy preservation as an algorithm. Instead, Camerala contributes a local-first browser-native implementation example that integrates remote task execution, face-derived feature extraction, task-event/frame-record timestamp alignment, operational landmark-availability checking, and local CSV/ZIP persistence without raw facial video retention for remote psychological/psychophysiological task experiments.''

\textbf{Location.}
Related Work; Contributions; System Architecture; References.

\condition{6}{Unclear claim of robustness against rigid head motion}
\textbf{Response.}
Thank you for this comment.
Unsupported robustness claims were removed.

\textbf{Revision.}
We removed unsupported robustness claims.
The revised manuscript does not claim that Camerala corrects motion-induced rPPG degradation.
Motion estimates are reported only as exported sensor-derived records for descriptive inspection.

\textbf{Key revised text.}
``This feature is used as a motion-related quality and behavior indicator. It is not an algorithmic correction for motion artifacts in rPPG.''

\textbf{Location.}
System Architecture, Feature Extraction Logic; Discussion.

\condition{7}{Insufficient description of the concrete usage scenario}
\textbf{Response.}
Thank you for the comment.
The intended usage scenario is now stated at the beginning of the paper and again in the Methodology.

\textbf{Revision.}
The manuscript now defines the target use case as a remote psychological or psychophysiological task experiment conducted while the participant is seated in front of a laptop, views task stimuli, and responds through the browser, with synchronized browser-based logging of behavioral events and webcam-derived features.

\textbf{Key revised text.}
``The target use case is a remote psychological or psychophysiological task experiment conducted while the participant is seated in front of a laptop.''

\textbf{Location.}
Introduction; Methodology, Target Usage Scenario and Exploratory Deployment.

\condition{8}{Definition of ``in-the-wild'' and ``naturalistic physiological sensing'' are not clear}
\textbf{Response.}
Thank you for this point.
The manuscript was revised to avoid giving the impression that Camerala was evaluated for unconstrained daily-life physiological sensing.

\textbf{Revision.}
Unqualified ``in-the-wild'' and ``naturalistic'' expressions were removed from the manuscript.
The revised manuscript describes the study as a remote psychological or psychophysiological task experiment conducted in laboratory and home settings.
It also explicitly states that the deployment is not continuous daily-life sensing: the participant was seated in front of a laptop, viewed task stimuli, and responded through the browser.

\textbf{Key revised text.}
``This use case is not unconstrained daily-life sensing: the participant is seated in front of a laptop, views task stimuli, and responds through the browser.''

\textbf{Location.}
Abstract; Introduction; Methodology; Discussion; Conclusion.

\condition{9}{The description of the system hardware is insufficient for reproducibility}
\textbf{Response.}
Thank you for this comment.
A hardware and software configuration table was added.

\textbf{Revision.}
The hardware/software table was revised using author-confirmed information: PC model, CPU, GPU, RAM, storage, OS, browsers, built-in webcam, requested camera constraints, browser-negotiated capture behavior, display resolution/refresh rate, fixed screen brightness, viewing distance, and default camera-control behavior.

\textbf{Key revised text.}
``Camera capture behavior: Browser-negotiated capture under the requested constraints; no manual exposure or white-balance constraints were applied.''

\textbf{Location.}
Methodology, Hardware and Software Configuration; Table 2.

\condition{10}{The description of the experimental environments is insufficient}
\textbf{Response.}
Thank you for the detailed request.
An environment and lighting conditions table was added.

\textbf{Revision.}
We limited the environmental description to setup/check information and camera-derived records.
The brightness range is not reported as lux or externally measured illuminance.
The revised section reports the laboratory/research room and private home room, daytime sessions, windows present, curtains closed, general artificial indoor lighting, fixed screen brightness, and camera-derived brightness ranges observed during setup/checks.

\textbf{Key revised text.}
``Because the deployment used camera-derived setup checks rather than external light-meter measurements, this paper reports brightness ranges only as author-confirmed setup/check information.''

\textbf{Location.}
Methodology, Deployment Environment and Lighting Conditions; Table 3.

\condition{11}{The validity of the cognitive framing manipulation is questionable}
\textbf{Response.}
Thank you for this important point.
The manuscript no longer interprets the negative, neutral, and positive instruction blocks as inducing evaluation pressure, reward motivation, cognitive effects, motivational effects, or psychological effects.

\textbf{Revision.}
The terminology was changed from cognitive framing manipulation to demonstration instructional task contexts.
The participant/autoethnography section now states that the author-participant knew the instructions were experimental.

\textbf{Key revised text.}
``They are used only as different instructional task contexts for demonstrating synchronized logging across task conditions.''

\textbf{Location.}
Participant section; Experimental Procedure and Instructional Task Contexts; Results.

\condition{12}{The central claim in the Discussion should be reconsidered}
\textbf{Response.}
Thank you for this comment.
The Discussion was rewritten around what the system demonstrated under the tested deployment and what remains unresolved.

\textbf{Revision.}
The revised Discussion focuses on the fact that the system was built and ran under laboratory and home settings, while the observed differences remain deployment-context differences that cannot isolate psychological or physiological effects.
Future validation requirements and design safeguards were added.

\textbf{Key revised text.}
``The results do not establish general deployability, robustness, or physiological/cognitive effects. Instead, they motivate design requirements and safeguards for future local-first deployments.''

\textbf{Location.}
Discussion.

\condition{13}{Several claims are not sufficiently supported and should be weakened}
\textbf{Response.}
Thank you for this guidance.
The Abstract, Results, Discussion, Limitations, and Conclusion were revised to substantially weaken unsupported claims.

\textbf{Revision.}
We rechecked the logs and analysis description and limited the revised results to indicators directly recomputable from the exported CSV files.
The manuscript no longer states that the study ``establishes'' a design pattern.
We replaced the ambiguous confidence-threshold wording with an implementation-faithful description of \texttt{mean\_quality} as an exported-window-level landmark-availability ratio.
The \texttt{quality} field is defined as a binary landmark-availability flag, and \texttt{mean\_quality} is the average of this flag within each exported 10-s window.
The current runtime forwards results to the logging callback only when FaceMesh landmarks are returned; therefore, the saved frame records in this deployment represent landmark-returned exported frames.
Thus, \texttt{mean\_quality} $\geq$ 0.8 means that landmarks were available in at least 80\% of exported frame records in that window.
This value is not a continuous MediaPipe confidence score, an indicator of rPPG measurement quality, or evidence of physiological measurement accuracy.
In the present deployment, all 299 exported windows had \texttt{mean\_quality} = 1.0; therefore, no window was excluded by this criterion.
The RT-based retained-trial count was updated to 541/600 after applying the 100--1500 ms reaction-time criterion.
The revised manuscript states that 600 trials were completed in total and that 541/600 trials were retained after the reaction-time criterion.
The previous exposure-related term was replaced with a frame-to-frame green-channel brightness-change indicator.
The manuscript also clarifies that this indicator is not a measure of camera exposure settings, external illuminance, environmental events, psychological or physiological state, or rPPG accuracy.

\textbf{Key revised text.}
``For each exported 10-s window, \texttt{mean\_quality} was computed as the mean of the \texttt{quality} flag over the exported frame records in that window. Thus, \texttt{mean\_quality} describes landmark availability within exported frame/window records and not face-detection availability across all camera frames.''

\textbf{Location.}
Abstract; Results; Discussion; Limitations; Conclusion.

\section*{Reviewer 1 Other Comments}
\textbf{Response.}
Thank you for the additional comments on figure and table readability.

\textbf{Revision.}
Table 1 was reformatted with clearer column widths.
The previous longitudinal and task-result figures contained overstrong embedded labels, so they were removed from the revised manuscript rather than retained.
The remaining results presentation relies on conservative text and tables without p-values, significance marks, or causal wording.

\textbf{Location.}
Table 1; Results; removed previous Figure 2, Figure 3, and Figure 4 from the revised manuscript.

\end{document}
